\section{The Blessedness Christ Teaches}

The Beatitudes give the desire for heaven a recognizable form. At the beginning of the Sermon on the Mount, Jesus pronounces blessed those whose condition can appear least secure: the poor in spirit, the meek, those who mourn, those who hunger for justice, the merciful, the clean of heart, the peacemakers, and those persecuted for righteousness (Mt 5:1--12). His promises include the kingdom of heaven, consolation, divine sonship, and the sight of God. The people he names must still wait for consolation, justice, and peace. Jesus calls them blessed in their present need because he promises their fulfillment in God. They can persevere in seeking those goods before they possess them.

Poverty of spirit places us before God as receivers. We possess nothing by which we could make him our debtor or secure his kingdom as an acquisition. Even a long life of fidelity remains sustained by gifts received. This poverty can be lived amid different outward circumstances, but it cannot remain indifferent to the use of material goods. The heart that knows its dependence upon God must learn to hold possessions in a way that serves his purposes. Wealth loses its power to confer final security; deprivation loses its power to prove that a person is forgotten. Receiving the kingdom requires trust in the Father who gives it. Examining the use of our possessions can show where we are seeking a security they cannot provide.

The promise to the meek has a similar depth. Meekness is strength brought under the government of charity. We need it whenever a good could be pursued through domination, humiliation, or revenge. The meek can defend the vulnerable and speak firmly because gentleness does not require the surrender of truth. What they relinquish is the claim to possess others by force of personality or power. Christ's promise of an inheritance exposes the insecurity beneath our grasping. We receive our lasting good from the Father. In preparation for that inheritance, we can learn to exercise authority as a service, to leave room for another person's response, and to endure a limit without making someone else pay for it.

Those who mourn are also called blessed. The promise does not require us to regard every sorrow as a sign of virtue; it directs sorrow toward the consolation of God. We mourn death, repent of sin, and grieve the injuries inflicted upon the innocent. Such grief can deepen our knowledge that the world needs a redemption greater than our own capacity to repair it. The God who consoles will bring his work to completion. Meanwhile, we can remain beside someone whose sorrow we cannot relieve, allowing him to grieve and trusting the consolation God promises.

Hunger and thirst for justice turn expectation into a sustained appetite for the good. We know what bodily hunger does: it gathers attention, makes the need concrete, and refuses to treat nourishment as an optional ornament. Christ uses that urgency to describe the desire for righteousness. We are to want the right ordering of our own lives and the good of those whom injustice harms. Satisfaction will come from God; the hunger he blesses nevertheless asks for our present action. A person who longs for heaven cannot be content that a weaker person is cheated, a family abandoned, or the truth traded for advantage. Our powers are limited, but their limitation does not cancel the work they can do. We can continue working where a wrong remains, asking God for the perseverance and help the task requires.

The merciful are promised mercy. In their conduct the gift received begins to bear fruit: another person's misery becomes a summons to help. Mercy can take the form of food, protection, instruction, pardon, or time freely given. Its manner must fit the actual need. Sometimes kindness requires a patient explanation; sometimes it requires protecting a person from the consequences of another's wrongdoing. The heavenly promise draws these acts into the generosity of God without turning them into a purchase price. We need mercy ourselves, and grace makes us capable of becoming its servants. Each genuine work of mercy prepares us for a communion in which another's good is welcomed as our own joy.

The clean of heart shall see God. Here Christ joins the highest promise to an inward transformation. Purity includes chastity, but the heart also needs cleansing from deceit, divided allegiance, and the wish to appear devoted while keeping its real treasure elsewhere. To desire the vision of God is to desire truth in the place from which our choices arise. We can let the Gospel uncover the intention beneath an outwardly admirable act. We can acknowledge that a service has become a means of influence or that religious knowledge has become a possession we use to rank people. Such discoveries are invitations to cleansing. God prepares the heart to receive himself by freeing it from the falsehoods with which it tries to secure a lesser good.

Peacemakers bear the family likeness of the God whose children they will be called. The peace they seek has room for truth, repentance, and the restoration of right relations. An appearance of calm sustained by intimidation cannot satisfy that desire. Nor can peace be built by endlessly preserving the quarrels through which we defend an injured importance. To make peace may require bringing a concealed harm into the light, asking pardon, or accepting the slow labor of rebuilding trust. The final communion gives this work its direction. We learn to want people reconciled in the good, even when reconciliation gives us no public credit and requires others to receive something we cannot supply.

The last Beatitude joins the promise of the kingdom to persecution for righteousness and fidelity to Christ. Opposition by itself does not prove that we are faithful; we can provoke it through our own injustice or folly. The Lord speaks of suffering borne for his sake. When fidelity truly costs us acceptance, livelihood, or safety, the promised kingdom declares that the cost cannot measure the worth of what is being given. Our life is held within a belonging deeper than the verdict of those who reject us. This is a hope the martyrs lived with their whole bodies. It can also sustain the quieter courage required to remain honest, protect an endangered person, or confess the faith when doing so carries a real loss.

The Beatitudes describe related changes in a person's life. Recognizing our own need for mercy prepares us to show it to another person; seeking truth in our own heart helps us work more honestly for peace. Their fulfillment is the Lord's gift, and their present beginnings depend upon his grace. Jesus follows them with the calling to be salt and light, so that human works may lead others to glorify the Father (Mt 5:13--16). Heavenly hope becomes visible through a life whose good points beyond itself. In an annual meditation, we can dwell upon the Beatitude that most closely touches our present responsibility and ask to live it more fully. The responsibility remains when we finish the meditation and calls for the grace we have sought in prayer.

